\chapter{Chapter 5}\label{ch:Chapter5}
\section{Protocolli End-To-End}
I protocolli end-to-end sono dei protocolli a livello trasporto e riguardano il \textit{Simple Demultiplexing} (UDP) e il \textit{Reliable Byte Stream} (TCP). Nel livello 3 si è visto come poter inviare pacchetti da un host a un altro situati in punti qualsiasi di una rete molto grande (astraendosi da dove e come essi siano collegati fisicamente), tuttavia non si è approfondito come evitare di perdere pacchetti, come riordinarli, come inviare duplicati del messaggio, limitare le dimensioni dei messaggi, ecc. Per cui al livello 4 si cerca di gestire queste lacune.
\section{Livello Trasport}
In questo livello di sviluppano algoritmi che trasformino le proprietà desiderabili della rete sottostante nell'alto livello di servizio richiesto dai programmi applicativi.Alcune funzionalità (come il QoS o il controllo della congestione) non si adattano ad un singolo livello, ma richiedono la collaborazione tra livelli.In alcuni casi, invece, non si ha un solo programma che gira, per cui si deve gestire anche un carico di lavoro maggiore e diversificato.In definitiva: tra i compiti principali del livello trasporto troviamo la pacchettizzazione,l'indirizzamento,l'affidabilità e il controllo sulla connessione.
\subsection{Indirizzamento end-to-end TCP/IP}
Il modo per verificare e identificare i processi a livello di Internet è tramite la \textbf{porta}.Il numero di porta è un numero compreso fra 0 e 65535 (16 bit) e ogni porta con tale numero viene considerata porta fisica, ovvero dispositivi fisicamente collegati ad una macchina.In ogni connessione sono coinvolti due numeri di porta, uno per ogni processo sui due host di comunicazione oltre agli indirizzi degli host su cui i processi sono in esecuzione.Se quindi si hanno molti processi su un host saranno necessare molte porte sullo stesso indirizzo di rete. \\
Per identificare univocamente un processo un'idea,anche se poco furba, può essere quella di fornire al numero di porta lo stesso numero del processo.Qui sorge un problema dato che se si cerca di fare in questo modo bisognerebbe che tutti i processi del mondo siano d'accordo sulla numerazione dei processi dunque è infattibile.
\subsection{Socket}
Un endpoint di comunicazione tra due host è chiamato Socket ed è creato tramite opportune chiamate di sistema.La socket è coppia di dati formata rispettivamente da un indirizzo ip e da una porta (\textit{indirizzo host:numero porta}). L'indirizzo ip identifica l'host mentre il numero di porta identifica il processo all'interno dell'host.Due processi, per comunicare, devono avere ciascuno un socket e utilizzare lo stesso protocollo di trasporto. In ogni istante, una comunicazione è \textbf{completamente identificata} dalla tupla: $<Ip_{A},port,Ip_{B},protocol>$.
\subsection{Modello Client/Server}
Spesso due processi comunicano tra loro tramite il modello Client/Server.
\begin{itemize}
  \item Server: è un processo che offre un servizio (es.:l'accesso ad una risorsa), è sempre in esecuzione ed in attesa delle richieste del client
  \item Client: è un processo che desidera utilizzare un servizio, può essere attivo solo quando necessario e invia le richieste al server.
\end{itemize}
\textbf{NB.:Client e Server sono processi e non host!} \\
Per iniziare la comunicazione il client deve conoscere la porta del server mentre il server scopre l'indirizzo del client quando viene contattato.Molto spesso il server utilizza delle porte note e comuni, come per esempio la porta 80 (utilizzata dai WebServer), mentre il client utilizza delle porte dinamiche scelte in modo casuale e possibilmente una per ciascuna comunicazione. Proprio per questa ragione la IANA ha stabilito degli intervalli per le porte:
\begin{itemize}
  \item WKP (Well Known Port): utilizzate dai server ufficiali. Possono essere utilizzate da processi con diritti di amministratore. (0..1023)
  \item Porte registrate: possono avere un uso standard ma non limitato ai processi di amministrazione (es.:i DBMS utilizzano queste porte, come la 3306 di MySQL). (1024..49151)
  \item Dinamiche: libere per tutti, solitamente assegnate dal kernel ai client che non specificano alcuna porta (49152..65535)
\end{itemize}
\subsection{Tipi di comunicazione}
Esistono diversi tipi di comunicazione implementabili con le socket, tra le principali troviamo:
\begin{itemize}
  \item Connection Oriented: è necessario stabilire una connessione prima di scambiare i dati, il che è molto costoso, in quanto si deve assicurare che durante tutta la comunicazione ci sia uno stato ricco delle informazioni.I datagrammi appartengono ad un'unica connessione, quindi sono più controllati e di conseguenza più affidabili (ne risente il setup).Nel modello TCP/IP questo viene implementato da TCP.
  \item Connectionless: in questo caso si va a replicare con il minimo sforzo il modello della rete di livello 3 sottostante, ogni messaggio (datagramma) viene inviato senza l'apertura di una connessione esplicita ottenendo quindi una risposta più rapida con lo svantaggio di una comunicazione meno affidabile. Nel modello TCP/IP questo viene implementato da UDP.
\end{itemize}
Altri tipi di comunicazione implementati (non completamente) sono:
\begin{itemize}
  \item Reliable Connectionless: è teoricamente possibile ma raramente implementato.Può essere simulato utilizzando TCP o aggiugendo controlli a UDP.I pacchetti sono affidabili e controllati. Nel modello TCP/IP è implementato tramite RUPD e RDP.
  \item Request/Reply: ogni richiesta del client è seguita da un corrispondente risposta del server (es.:richiesta HTTP da server e risposta). Ogni richiesta/risposta è indipendente dalle altre e la richiesta può essere rifiutata.E' una tecnica molto diffusa nelle architetture di comunicazioni attuali.Nel modello TCP/IP è implementato nel protocolli a livello applicazione (es.:RPC,RMI,SOAP,REST,ecc.)
\end{itemize}
\section{UDP}
Il protocollo UDP (User Datagram Protocol) è definito anche come demultiplexer semplice.Esso estende il servizio di consegna da host a host,della rete sottostante, in un servizio di comunicazione da processo a processo e aggiunge un livello di demultiplexing che consente a più processi applicativi, su ciascun host, di condividere la rete.Non vi è alcuna garanzia oltre a quella di IP \\
Il formato del datagramma UDP è composta da un header e un payload. L'header è composto da Source port number (16bit), Destination port number (16bit), Total lenght (16bit) e Checksum (16bit) per un totale quindi di 8Byte.
\subsection{Calcolo Checksum}
Il checksum se presente (quindi non nullo),è calcolato sommando \textit{header UDP + dati + intestazione pseudo-IP}.L'intestazione psudo-IP è un riassunto dell'intestazione IP vera e propria ma contenente solo: indirizzi IP,protocollo (per UDP è il 17),lunghezza totale del datagramma IP, mentre altri dati sono omessi (es.:frammentazione,TTL,ecc.). Esso identifica la comunicazione ovvero la quintupla indirizzo ip + porta.Questo header in realtà non viene trasmesso ma bensì viene utilizzato soltanto per il calcolo del checksum.Nella realtà viene trasmesso un header di 20byte e se per caso durante la connessione l'indirizzo IP dovesse cambiare, si dovrebbe ricalcolare il checksum.
\subsection{Api per i servizi UDP}
Supponiamo che un pacchetto UDP venga accettato.Esso viene inserito in appositi buffer, è possibile infatti immaginare il modello UDP come un insieme di code per datagrammi.I buffer utilizzati sono uno di ingresso e uno di uscita in modo da poter far arrivare i datagrammi,processarli e poterne costruire uno nuovo per rispondere o inoltrare. Le code sono accessibili dalle applicazioni utilizzand due system call: \textit{sendmsg()} e \textit{recvmsg()}.Quando si esegue la recvmsg(), se nella coda di input è presente un datagramma questo viene spostato su buffer interni per poi essere processato, se vuota allora l'host si blocca in attesa di un nuovo datagramma.Dato che recvmsg() è una chiamata bloccante entra in gioco il concetto di sincronizzazione tra client e server in particolare su quanto è fondamentate l'utilizzo di sendmsg(). Quest'ultima permette di soddisfare le operazioni senza ostacolare altre funzioni quando serve inviare un nuovo datagramma.
\subsection{Demultiplexer semplice}
I datagrammi, quando giungono all'host, vengono controllati e inseriti nella coda associata alla porta indicata. Siccome non si sa quanto sono grandi le code potrebbe accadere che un pacchetto venga scartato silenziosamente perchè la dimensione della coda è limitata e già piena con altri datagrammi.Proprio per questi motivi potrebbe succedere che un pacchetto non arrivi a destinazione
\begin{example}
  Invio di un video/audio attraverso UDP.Essendo senza controlli di flusso se il buffer si satura perchè la sorgente continua ad inviare pacchetti questi verrebbero scartati e nuovamente richiesti dal destinatario.
\end{example}
Altro punto importante, i datagrammi non sono numerati e non appartengono a una sessione/connessione per cui sono consegnati in modo indipendente.Non vi è alcun controllo di errore di conseguenza il mittente non viene avvisato se vengono persi dei pacchetti oppure il destinatario è congestionato.Come se non bastasse vi è solo incapsulamento IP e potrebbe verificarsi frammentazione.
\subsubsection{Riassumento..}
UDP è orientato alle transazioni (transaction-oriented), adatto a semplici protocolli di richiesta come il DNS o l'NTP,è semplice,stateless,buono per il boostrap o per altri scopi senza un stack di protocollo completo (es.:DHCP e TFTP) e adatto a un numero molto elevato di client, come nelle applicazioni di streaming multimediale dove non è un problema perdere qualche pacchetto ma i dati devono arrivare velocemente.E' adatto a comunicazioni unidirezionali e per l'implementazione di specifici protocolli di trasporto end-to-end nello spazio utente.
\section{TCP}
Implementato nel 1974 dalla IEEE, \textit{Transfer Control Protocol} offre servizi su flusso di byte affidabili (reliable byte stream) e orientamento alla connessione.Il funzionamento prevede la creazione di una sorta di imbuto o pipe tra un processo e un altro per instaurare una connessione.Questa astrazione garantisce che i dati vengano tutti inviati in ordine e senza ripetizioni.TCP gestisce da se tutte le problematiche relative al recupero di eventuali dati andati persi e di ordinamento dei segmenti dei dati.la comunicazione è bidirezionale dato che una socket ha due pipe una di andata e una di ritorno.
\subsection{Flow control vs Congestion control}
TCP si occupa di:
\begin{itemize}
  \item Controllo del flusso:impedisce al mittente di superare la capacità del ricevitore, il controllo viene delegato e gestito dai livelli superiori (questo c'era gia nella Sliding Window a livello 2).
  \item Controllo di congestione:consiste nell'impedire che una quantità eccessiva di dati venga immessa in rete, causando così un sovraccarico degli switch o dei link.La congestione è quando un router viene raggiunto da diversi flussi di pacchetti che non riesce gestire,la conseguenza é quella di scartare i pacchetti in ingresso.Tutto questo è all'oscuro del mittente e destinatario perchè avviene nei router (che si trovano in mezzo).
\end{itemize}
Il TCP deve affrontare entrambi i problemi e di conseguenza, la larghezza di banda e i ritardi end-to-end sono definiti dal destinatario e dal link/switch più lento lungo il percorso dalla sorgente alla destinazione. Il controllo di flusso è relativamente facile da gestire,basta implementare un meccanismo di avviso di riempimento imminente nella coda da parte del destinatario verso il mittente.La gestione della congestione invece è molto più complessa, in quanto implica conoscere lo stato della rete quando nella realtà non lo si conosce (non si sà quali sono i router intermedi).
\subsection{Problemi end-to-end}
Il cuore del TCP è l'algoritmo \textit{Sliding Window}, i cui problemi da affrontare sono dati dal fatto che TCP funziona su Internet piuttosto che su un collegamente punto-punto; TCP infatti supporta connessione logiche tra processi in esecuzione su due computer diversi in Internet che possono avere tempi di RTT molto diversi e pacchetti disordinati.Per il controllo del flusso, il TCP ha bisogno di un meccanismo che permetta a ciascun lato di una connessione di sapere quali risorse l’ altro lato è in grado di applicare alla connessione (upgrade della finestra scorrevole). Mentre, per il controllo della congestione, il TCP ha bisogno di un meccanismo che consenta al lato mittente di conoscere la capacità della rete.
\subsection{Segmento TCP}
Il TCP è un protocollo orientato ai byte, il che significa che il mittente scrive i byte in una connessione TCP e il destinatario legge byte dalla connessione TCP (si verifica perdita di struttura poichè i dati vengono presi dal buffer e trattati come un fluido. Anche se vengono trasformati in byte non importa in che quantità o in che dimensione di blocchi arrivano, essi vengono ricomposti com'erano).Inoltre non vengono mantenuti i confini dei messaggi, cosa che gli altri protocolli fanno.Sebbene il flusso di byte descriva il servizio che il TCP offre ai processi applicativi,il TCP non trasmette di per sè singoli byte su Internet.Sull'host di origine il TCP bufferizza una quantità di byte,dal processo di invio, sufficiente a riempire un pacchetto di dimensioni ragionavoli e poi invia questo pacchetto al suo peer sull'host di destinazione.Poi, sull'host di destinazione, TCP svuota il contenuto del pacchetto in un buffer di ricezione dal quale,il processo ricevente, legge a suo piacimento i byte.I pacchetti scambiati tra i peer TCP sono chiamati \textit{segmenti} (payload) ovvero una fetta di byte. \\

Il protocollo TCP può essere visto come un modello produttore-consumatore, in cui si hanno dei buffer suddivi in sezioni ogniuno con delle caratteristiche diverse.Il buffer \textit{Sent} (in uscita) ha tre colorazioni: bianco per la parte libera,in cui l'applicazione può scrivere i dati, rosa per i dati scritti nel buffer e che non sono ancora stati spediti e grigio per i dati inviati e da cui si attende un riscontro (è un backup dato che il destinatario non ha ancora notificato il loro arrivo tramite ACK).Il buffer \textit{Not read} (in entrata), invece, ha solo due colorazioni:bianco per la parte libera e rosa per la parte di dati ricevuta e che non è stata ancora utilizzata.Per sapere a che punto si è arrivati con la lettura/scrittura e con la spedizione/arrivo si utilizzano dei puntatori.Il mittente prende una parte libera del proprio buffer e inizia a scriverci, man mano che che invia si segna in grigio quali byte non hanno ancora ricevuto un riscontro.Quando il riscontro positivo arriva, lo spazio del buffer occupato da questi byte viene liberato.Il destinatario a sua volta riempie i byte liberi del proprio buffer, notifica il mittente di cio che ha/non ha ricevuto e se ha ancora spazio per ricevere.A mano a mano che i buffer vengono riempiti/svuotati si spostano i puntatori.
\subsection{Hearder TCP}
Nell'header TCP, di 20byte (default), sono segnati gli indirizzi e le rispettive porte, seguiti da un numero di sequenza (dove devono essere collocati i dati del payload all'interno della sequenza), c'è il campo acknowledgment (ACK) per l'altra direzione, in modo da informare il destinatario fin dove il mittente è stato notificato della ricezione (tecnica chiamata piggybacking, in cui invece che inviare pacchetti separati della stessa cosa o di qualcosa di simile, li si mettono assieme in un unico pacchetto, in questo caso ACK e dati effettivi),l'HLEN (lunghezza dell'header, solitamente 5 word), c'è poi una serie di bit di controllo (flag SYN,FIN,RESET,PUSH,URG e ACK), il campo che indica la grandezza della Window (memoria disponibile del buffer di ricezione, per il controllo di flusso, altro piggybacking), il campo checksum (uguale a quello UDP) e altre informazioni di controllo.
\subsubsection{Nel dettaglio}
I flag SYN e FIN sono utilizzati rispettivamente per stabilire e terminare una connessione TCP, ACK è un flag che viene settato ogni volta che il campo \textit{Acknowledgment} è valido; implica che il destinatario deve prestarvi attenzione. Il flag URG indica che questo segmento contine dati urgent; se è settato il campo \textit{UrgPtr} indica dove iniziano i dati non urgenti contenuti nel segmento.Solitamente i dati urgenti sono contenuti nella parte anteriore del corpo del segmento e si estendono fino ad una valore di UrgPtr presente.Il flag PUSH indica che il mittente ha invocato l'operazione push,questo va a dire alla parte ricevente del protocollo TCP che deve avvisare il processo di ricezione di questo fatto.Il flag RESET indica che il destinatario è confuso, ha ricevuto un segmento che non si aspettava di ricevere e quindi vuole interrompere la connessione.Lo stato di confusione può essere determinato dal fatto che TCP è stateful e c'è bisogno di sincronizzare gli endpoint.
\subsection{Diagramma degli stati TCP}
Una comunicazione orientata alla connessione deve mantenere alcune informazioni sui dati,sui segmenti mancanti e molto altro. Ogni socket è in uno stato in cui le transizioni da diversi eventi, quali: l'esecuzione di comandi a livello applicazione sul socket, come \textit{accept()},\textit{connect()} e \textit{close()}, oppure la ricezione di segmenti con flag appropriati.Le transizioni a loro volta possono causare l'invio di segmenti con flag appropriati.Sempre ricordando che le socket  devono avere stati coerenti fra loro. \\
TCP può essere visto come un automa a stati, i cui stati (finiti) e le transizioni di stato di un socket TCP sono definiti dal protocollo per mezzo di un DFA (automa a stati finiti).Ogni transizione è innescata da un evento e può produrre la trasmissione di un segmento. Il tutto si articola in 3 specifiche fasi:
\begin{enumerate}
  \item \textcolor{cyan}{\underline{Handshaking}}:creazione della connesione
  \item \textcolor{lime}{\underline{Trasferimento dei dati}}:quando i dati fluiscono
  \item \textcolor{orange}{\underline{Chiusura}}:fine della connessione
\end{enumerate}
Una connessione può rimanere \textbf{ESTABLISHED} per un tempo illimitato (anche per sempre).
\subsubsection{Fase 1(handshake):}
Prima di inviare qualsiasi dato due processi devono stabilire una connessione.Un processo (\textit{responder} o server) attende la una connessione, eseguento una funzione TCP adatta per l'apertura passiva (in Unix è la accept()).L'altro processo (\textit{initiator} o client) invece inizia la connessione eseguendo l'apertura attiva(in Unix è la connect()).A questo punto, i due livelli TCP eseguono il three-way hankshake:
\begin{enumerate}
  \item il server attende in stato di LISTEN
  \item quando il client esegue l'apertura passiva invia il primo segmento (SYN) che ha un numero di sequenza iniziale spostandosi nello stato SYN-SENT
  \item il server risponde con SYN+ACK: riconosce il numero di sequenza del client e imposta il suo numero di sequenza (anche in questo caso nessun payload) cambiando stato in SYN-RECEIVED.
  \item il client risponde con ACK: riconosce il numero di sequenza del server e passa allo stato ESTABLISHED
  \item  quando il server riceve l'ACK passa anch'esso in ESTABLISHEDA,per cui la connessione è stabilita e i segmenti che trasportano il payload di dati reali possono essere inviati in entrambe le direzioni.
\end{enumerate}
Ci sono casi in cui entrambi gli endpoint eseguono un richiesta di connessione attiva.In questo caso uno dei due client, nel vedersi arrivare un segmento SYN, potrebbe decidere di rifiutare la richiesta per non cambiare la propria natura e quindi invia un RESET.L'altro endpoint accetta e passa dallo stato SYN-SENT a SYN-RECEIVED.I due endpoint si scambiano quindi i due segmenti ACK e passano entrambi allo stato ESTABLISHED.Questo rarissimo cambio di procedura prende il nome di \textit{Simultaneous Open}, dato che entrambi i client effettuano una open.In realtà, un client non sa se un altro endpoint (presumibilmente un server) è in apertura passiva fintanto che non cerca di connettersi e se non lo è riceve un segmento RESET.\\
I segmenti che vengono inviati presentano un numero random di sequenza.I segmenti ACK riportano lo stesso numero del campo sequenza dei SYN a cui fanno riferimento,ma incrementato di uno, in questo modo è possibile riconoscerli.
\subsubsection{Fase 2(trasferimento dei dati):}
Entrambe le parti sono in ESTABLISHED.Ogni segmento trasporta dati in una direzione e la conferma dei dati nella direzione opposta(dato che i flussi sono indipendenti).Prendiamo in esempio il numero sequenza di partenza 8001.Supponiamo che il pacchetto sia lungo 1Kbyte circa,quidni va dal byte 8001 al byte 9000 (il primo segmento che il client invia).Successivamente viene inviato un altro segmento che va dal byte 9001 al 10000.Entrambi hanno l'ACK con valore 15001(valore dell'handshake).A questo punto arriva la risposta dal server, che conferma la ricezione dei pacchetti tramite un ACK con numero di sequenza di 15001 e numero di ACk 10001(la lunghezza può variare dal byte 15001 al 17000).Questo segmento appena inviato contiene il nuemro cumulativo dei pacchetti fino a quel momento ricevuti.Quando il client riceve l'ACK sa che può iniziare a liberare i dati di backup e in contemporanea invia un ACK che conferma la ricezione della risposta precedente.Il flad PUSH, che è stato utilizzato dal client nei primi due segmenti, indica al server di notificare la ricezione e lo forza a tenerlo aggiornato su: sequenze di byte che riceve,l'invio di dati il più velocemente possibili senza creare backup ("data pushing", assenza di dati).\textbf{NB:l'invio multiplo di dati prima o poi viene interrotto se allo scadere del timeout non arriva un ACK che ne confermi la ricezione}. \\
Come si può notare, il TCP è strettamente sequenziale.A volte è utile consegnare i dati fuori sequenza;l'applicazione può chiedere una consegna urgente durante un invio,se viene accettata, il TCP del mittente colloca i dati urgenti all'inizio del segmento successivo,attivando il bit URG e impostando il puntatore urgente.Il TCP del destinatario trova URG=1,separa i dati urgenti da quelli standard e li consegna immediatamente, anche se sono fuori ordine.
\subsubsection{Fase 3(chiusura):}
Dopo un po' di tempo, uno dei due host può richiedere di chiudere la connessione.La chiusura è causata dalla system call $close$ su una delle due parti,che si chiamerà \textit{initiator}.L'altra parte risponde eseguendo la sequenza di risposta.Nel raro caso in cui entrambe le parti inizino contemporaneamente la sequenza initiator, si ha la chiusura simultanea.Un host(initiator) avvia la procedura di chiusura inviando un segmento con FIN=1 e passando allo stato FIN-WAIT-1.Questo segmento può contenere l'ultimo payload di dati dopo del quale nessun altro dato può essere inviato dall'host.L'altro host(responder) invece, conferma all'initiator la ricezione dell'ACK,va in stato CLOSE-WAIT e notifica la chiusura all'applicazione.L'initiator quindi si porta in FIN-WAIT-2.L'applicazione sul responder può ancora inviare dati all'altro host,che continuerà a confermarli.Per finire,l'appliczione sul responder chiude il socket,quindi quest'ultimo invia un FIN all'initiator e passa allo stato LAST-ACK.L'initiator riceve FIN,invia il suo ACK e si sposta in TIME-WAIT.Il responder chiude il socket quando riceve l'ACK. \\
Chi ha iniziato la chiusura e che,a questo punto,si trova in TIME-WAIT, ha un certo tempo prima di passare allo stato di CLOSED questo per evitare di tornare troppo presto nello stato iniziale.Il pericolo è che:
\begin{itemize}
  \item se l'host non passase per lo stato TIME-WAIT potrebbe essere che esso chiuda la connessione e poi magari il suo ACK vada perso,con la conseguenza che,dato che una volta chiusa la connessione gli host non ricordano più nulla di essa, si riapra una vecchia connessione di cui si sono perse delle informazioni
  \item potrebbero esserci ancora vecchi pacchetti in giro per la rete, di cui non è chiaro a quale connessione si riferisca si rischi quindi che interferiscano con altre conunicazione degli host che prima erano interessati.
\end{itemize}
Per cui si importa un tempo di attesa chiamato \textit{Maximum Segment Lifetime (MSL)}, che ha una durata di circa 1 o 2 minuti.
\subsection{Rivisitazione della sliding window}
La variante TCP dell'algoritmo della sliding window ha diversi scopi:
\begin{itemize}
  \item garantisce la consegna affidabile dei dati
  \item assicura che i dati siano consegnati in ordine
  \item impone il controllo del flusso tra il mittente e il destinatario
\end{itemize}
\textbf{Usare appunti ilenia + figura}
\subsection{Protezione contro il wraparound}
Come previsto dal protocollo Sliding Window, bisogna assicurarsi che non ci siano due dati in sospeso con lo stesso numero di sequenza.La condizione è: \textit{larghezza della finestra < numero di sequenza massimo / 2} equivalente a \textit{numero di sequenza massimo > 2 * laghezza della finestra}.In TCP si ha un \textit{SequenceNum} di 32 bit e una \textit{AdvertisedWindow} di 16 bit.Quindi, il requisito dell'algoritmo della sliding window è soddisfatto perchè $2^{32} >> 2*2^{16}$.Le recenti implementazioni di TCP hanno più bit per AdvertisedWindow e la condizione resta valida.Tuttavia il meccanismo di avere uno spazio di numeri di sequenza maggiore della window potrebbe non bastare nel caso di connessioni "riavvolgenti"(wraparound).Infatti,se si trasmette un flusso di byte fino all'esaurimento dello spazio e  si necessita di ricominciare la sequenza,bisogna stare attenti che la ripetizione dei vecchi numeri non sia ancora dentro al buffer e quindi benga interpretata come un reinvio dei vecchi pacchetti, che verrebbero scartati.Per risolvere questo inconveniente, si aggiung un timeout di 120 secondi, chiamato \textit{durata massima del segmento (MSL)}, ovver il tempo massimo che i pacchetti possono sopravvivere in Internet.Il risci di riavvolgimento del numero di sequenza dipende dalla velocità di trasmissione dei dati su Internet.Esistono dei casi in cui i delay sono maggiori di 120 secondi, per esempio la rete è molto grande oppure in base al mezzo trasmissivo utilizzato.Per questo motivo,TCP è fatto per funzionare con delay molto maggiori di quello appena citato.
\subsection{Mantenere la pipe piena}
Il campo AdvertisedWindow deve essere suffiientemente grande da consentire al mittente di mantenere la pipe piena (chiaramente il destinatario è libero di non aprire la finestra tanto quanto lo consente il campo AdvertisedWindow).Se il destinatario ha spazio sufficiente per il buffer, la finestra deve essere apera abbastanza per consentire un \textit{ritardo * larghezza di banda} carico di dati.
\begin{example}
  Un campo a 16 bit consente un finestra di 64 Kbyte, che potrebbe non essere sufficiente per le "linuee lunghe", utilizzando la formula definita sopra si otterrebbero $ritardo * lunghezza di banda = 12500 byte$. TCP consente di impostare un \textit{fattore di scala} (fino a $2^{14}$) durante il protocollo di handshake.Pertanto la finestra può arrivare a 1GB ($2^{16} * 2^{14} = 2^{30}$).
\end{example}
\subsection{Attivazione della trasmissione}
Quando si deve trasmettere un segmento entra in gioco l'astrazione di flussi di byte,supportata da TCP.I programmi applicativi scrivono i byte nei flussi e poi spetta al TCP decidere se ha abbastanza byte per inviare un segmento. \\
Il TCP dispone di tre meccanismi per attivare la trasmissione di un segmento:
\begin{enumerate}
  \item TCP mantiene una dimensione \textit{massima del segmento(MSS)} variabile e invia un segmento non appena ha raccolto i byte MSS dal processo di invio.La MSS è solitamente impostata sulla dimensione del segmento più grande che il TCP può inviare, senza causare la frammentazione dell'IP locale.\textit{MSS = MTU della rete direttamente connessa - (header TCP + header IP)}.
  \item il processo mittente ha chiesto esplicitamente a TCP di inviarlo(TCP supporta l'operazione di push).
  \item quando scatta il timer.Il segmento risultante contiene tanti byte quanti ne sono attualmente bufferizzati per la trasmissione(ovviamente fino al limite definito da MSS e EffectiveWindow).
\end{enumerate}
\subsection{Silly Window Syndrome (SWS)}
L'algoritmo TCP a SlidingWindow di base non prevede una dimensione minima per i segmenti trasmessi.La SWS è una situazione in cui vengono inviati molti segmenti piccoli e inefficienti(più overhead), piuttosto che pochi segmenti grandi.Può verificarsi quando un destinatario pubblicizza dimensioni di finestrra troppo piccole o quando un trasmettitore è troppo aggressivo nell'inviare immediatamente quantità molto piccole di dati.Questo non si tratta di un fallimento dell'algoritmo, ma di un'inefficienza dovuta all'overhead della rete.Il caso peggiore è quando: AdvertisedWindow è 0(trasmissione ferma per buffer pieno),il processo di ricezione consuma 1 byte,quindi il ricevitore invia al mittente che AdvertisedWindow è 1;il mittente vedendo che si è liberato un byte di spazio invia aggressivamente un segmento di 1 byte,chiudendo nuovamente la finestra e costringendo la trasmissione a fermarsi(il buffer del ricevente è saturo);il tutto poi si ripete.QUesto porta a una sequenza di segmenti da 1 byte(ciasucono con 40 byte di overhead di header).Sarebbe ideale se il mittende aspettasse un po',prima di inviare un nuovo segmento,cosi da permettere al processo ricevente di consumare byte,aprire la finestra e di conseguenza permettere al mittente di inviare un singolo segmento più grande. \\
Tuttavia questa attesa introduce un rallentamento che potrebbe risultare scomodo per applicazioni interattive come Telnet e SSH, in cui ci si aspetta una risposta immediata in seguito al lancio del comando.Serve quindi una via di mezzo per quelle applicazioni che possono inviare grosse quantità di dati a discapito di un rallentamento e per quelle che devono essere efficienti e veloci.
\subsection{Algoritmo di Nagle}
John Nagle ha introdotto un'elegante soluzione per stabilire la connessione.L'idea chiave è quella di inviare un ACK al mittente durante il transito dei dati,mediante TCP.Questo ACK potrà essere trattato come un timer che scatta, facendo così partire la trasmissione di altri dati.Qualche istante prima di inviare l'ACK si attente il tempo necessario per permettere ai buffer di svuotarsi,finito lo svuotamento il destinatario notifica il mittente che è libero e non appena l'ACK arriva si procede con l'invio effettivo di nuovi dati.
\subsection{Ritrasmissione adattiva e calcolo del timeout}
Un segmento deve essere ritrasmesso se non si riceve una conferma per un certo periodo di tempo,bisogna impostare il timeout.L'algoritmo della SlidingWindow (livello 2)calcola il timeout in base alla tecnologia utilizzata e alla distanza dei due endpoint,risultando un calcolo semplice e deterministico.Al livello 4 è più complesso e tende a cambiare nel tempo,quindi non è sempre prevedibile e dipende dal percorso dei pacchetti, che potrebbero anche andare persi. \\
L'algoritmo misura l'RTT campione per ogni coppia segmento/ACK e calcola la media ponderata del RTT,ogni volta che un pacchetto viene inviato si tiene conto di quanto tempo server per avere un riscontro.Si ricalcola quindi, volta per volta la media,per cui si ha:
\begin{equation}
  EstRTT := \alpha * EstRTT + (1 - \alpha) * SampleRTT,\: con\: \alpha \: tra \: 0,8 \: e \: 0,9
\end{equation}
A questo punto si imposta il timeout in base al balore di \textit{EstRTT} ovvero $TimeOut = 2 * EstRTT$.Si nota che la formula è ricorsiva e necessita di un valore EstRTT iniziale(quello campionato durante l'handshake).\\
Questo algoritmo,però,presenta due problemi:
\begin{itemize}
  \item L'ACK non conferma realmente una trasmissione ma solo la ricezione dei dati
  \item è impossibili decidere a quale trasmissione l'ACK deve essere associato quando viene ritrasmesso un segmento.
\end{itemize}
\subsubsection{Algoritmo di Karn-Partridge:}
Una soluzione è quella di utilizzare l'algoritmo di Karn-Partridge(obbligatorio in TCP),che dice di non campionare l'RTT durante la ritrasmissione,ma di \underline{raddoppiare il timeout dopo ogni ritrasmissione}.L'algoritmo ignora i segmenti ritrasmessi quando aggiorna la stima del tempo di andata e ritorno(RTT).La stima del RTT si basa solo sui riscontri non ambigui,ovvero quelli per segmenti inviati una sola volta. \\
Il raddoppio del tempo serve in quelle situazioni in cui TCP invia un segmento dopo un forte aumento del ritardo.Se si utilizzasse la stima del RTT precedente il TCP dovrebbe calcolare un timenout che potrebbe ignorare l'RTT di tutti i pacchetti ritrasmessi; quindi la stima dell'RTT non verrà mai aggiornata e TCp continuerà a ritrasmettere ogni segmento,senza mai adattarsi all'aumento del ritardo.La soluzione a questo problema consiste nell'incorporare i timeout di trasmissione con una strategia di backoff a timer,che calcola un timeout iniziale,se il timer scale e causa una ritrasmissione, il TCP aumenta il timeout generalmente di un fattore due.Questo algoritmo è estremamente efficace nel bilanciare le prestazione e l'efficienza nelle reti con  un'elevata perdita di pacchetti,anche se idealmente non sarebbe necessario:le reti che presentano tempi di andata e ritorno elevati e timeout di ritrasmissione dovrebbero essere analizzate con tecniche di analisi delle cause principali.
\subsection{Presationi del TCP}
Per quando riguarda le prestazioni del TCP si tiene contro di due metriche:latenza e throughput.
\begin{example}
  Se si avesse una configurazione sperimentale con Zeon 2,4 GHz,una Gigabit Ethernet in aggregazione di link(Linux) e un limite del canale datalink di quasi 2Gbps full duplex,senza perdite,si avrebbe che:più grande è l'MSS più alto è il throutputma sopra 1KB non aumenta in modo sostanziale.L'MSS tipico è di 1460 byte(adatto quando l'interfaccia ha MTU=1500) che si mantiene al di sotto di 2Gbps(bottleneck).Nel complesso TCO può raggiungere un elevato throughput dove a volte il bottleneck è la memoria stessa.
\end{example}
